% (C) 2026 Ralf Stephan, in collaboration with Claude Code.
% Released under CC0 1.0 Universal (public-domain dedication).
% See https://creativecommons.org/publicdomain/zero/1.0/

\documentclass[11pt]{article}
\usepackage[T1]{fontenc}
\usepackage[utf8]{inputenc}
\usepackage{lmodern}
\usepackage[margin=2.5cm]{geometry}
\usepackage{amsmath,amssymb}
\usepackage{booktabs,tabularx}
\usepackage{braket}
\usepackage[hidelinks]{hyperref}

% Lean names and paths: verbatim, breakable at . _ /
\DeclareUrlCommand\lean{\urlstyle{tt}}
\def\UrlBreaks{\do\.\do\_\do\/}
\newcolumntype{L}{>{\raggedright\arraybackslash}X}
\newcolumntype{B}{>{\raggedright\arraybackslash}p{0.42\textwidth}}
\emergencystretch=2em

\title{Lean formalization of an Assumptions of Physics technical brief:\\ log}
\author{}
\date{}

\begin{document}
\maketitle

\noindent
This log records which statements of the technical brief in this directory
(\url{https://assumptionsofphysics.org/resources/briefs}) were checked in Lean, where
the formal version differs from the brief, and any problems found in the brief.

\section*{Setup}

\begin{itemize}
\item Lean \texttt{v4.36.0-rc1}; Mathlib at commit \texttt{2cabef97}; physlib (\texttt{master})
  and TauCeti (commit \texttt{eb4dd24}, on the same Mathlib) as further dependencies.
\item Build: \texttt{lake build} in the project root.
\item Every file is checked with \texttt{\#print axioms} on its main results. A result is
  recorded here only if it depends on no axioms beyond \lean{propext},
  \lean{Classical.choice} and \lean{Quot.sound}, and so has no \lean{sorry}.
\item Imports are kept to the Mathlib modules actually needed.
\item Material that belongs in Mathlib is kept in \lean{ForMathlib/}, in Mathlib's directory
  layout (library \lean{ForMathlib}), and imported by the brief files.
\end{itemize}


%======================================================================

\section{Superpositions correspond to multiple decompositions}

\paragraph{Files.}\leavevmode
\begin{itemize}
\item File: \lean{briefs/SuperpositionAndMultipleMixtures/SuperpositionAndMultipleMixtures.lean}
\item Library: \lean{SuperpositionAndMultipleMixtures}
\item Imports: \lean{ForPhyslib.QuantumMechanics.Wigner.Density} (for \lean{proj},
  $\ket\psi\bra\psi$, and its parallelogram law),\newline
  \lean{ForMathlib.Analysis.InnerProductSpace.Projection.FiniteDimensional}
\item New \lean{ForMathlib} material:\newline
  \lean{ForMathlib/Analysis/InnerProductSpace/Projection/FiniteDimensional.lean}
  (\lean{Submodule.mem_span_of_forall_inner_eq_zero}: a vector orthogonal to everything that is
  orthogonal to a finite set lies in its span).
\end{itemize}

\paragraph{Setting.} $E$ is a complex inner product space of any dimension, not necessarily
complete. For a vector $\psi$ the pure state is \lean{pureState}~$\psi = \rho_\psi =
\ket\psi\bra\psi/\braket{\psi|\psi}$ (it is $0$ at $\psi = 0$), and \lean{mixture}~$p\ \psi\ q\
\chi = p\rho_\psi + q\rho_\chi$. An operator is mixed (\lean{IsMixed}) if it is not
$\rho_\psi$ for any $\psi \ne 0$. The conditions are:
\begin{itemize}
\item \lean{QSupPur}~$\phi\ \psi_1\ \psi_2$: $\psi_1, \psi_2$ linearly independent and
  $\phi = c_1\psi_1 + c_2\psi_2$ with $c_1, c_2 \ne 0$;
\item \lean{QSupMix}~$\phi\ \psi_1\ \psi_2$: there are $p_1, p_2, p_\phi, p_{\hat\phi} \in (0,1)$
  with $p_1 + p_2 = p_\phi + p_{\hat\phi} = 1$ and a state $\hat\phi \ne 0$ such that
  $\rho = p_1\rho_{\psi_1} + p_2\rho_{\psi_2}$ is mixed, $\rho = p_\phi\rho_\phi +
  p_{\hat\phi}\rho_{\hat\phi}$, and $\rho_\phi \notin \{\rho_{\psi_1}, \rho_{\psi_2}\}$.
\end{itemize}

\paragraph{Results.}
\begin{center}
\begin{tabularx}{\textwidth}{@{}BL@{}}
\toprule
Brief & Lean \\
\midrule
eq.~(2), $\ket\phi\bra\phi + \ket{\hat\phi}\bra{\hat\phi} = 2|c_1|^2\ket{\psi_1}\bra{\psi_1} +
  2|c_2|^2\ket{\psi_2}\bra{\psi_2}$ & \lean{proj_add_proj_hat} \\
eqs.~(3), (4), QSUP-PUR $\Rightarrow$ QSUP-MIX & \lean{QSupPur.qSupMix}, with the weights of
  eq.~(3) \\
``it is supported by a two dimensional subspace'' & \lean{mem_span_of_mixture_eq}: if
  $q\rho_\phi + q'\rho_{\phi'} = p_1\rho_{\psi_1} + p_2\rho_{\psi_2}$ with $q > 0$, $q' \ge 0$,
  then $\phi \in \operatorname{span}\{\psi_1, \psi_2\}$ \\
QSUP-MIX $\Rightarrow$ QSUP-PUR & \lean{QSupMix.qSupPur}, for $\phi \ne 0$ \\
Theorem~1 & \lean{qSupPur_iff_qSupMix}, for $\phi \ne 0$ \\
(used for QSUP-MIX) & \lean{isMixed_mixture}: a mixture of two linearly independent states with
  positive weights is not pure \\
\bottomrule
\end{tabularx}
\end{center}

\paragraph{Differences from the brief.}
\begin{itemize}
\item The equivalence is proved for each triple $\phi, \psi_1, \psi_2$ with $\phi \ne 0$; no
  other vector needs to be nonzero, and $E$ need not be complete or finite-dimensional.
\item The brief writes $\rho$ as a density operator. The formalization does not use the trace:
  $\rho$ is a convex combination of two pure states, which is all the argument needs.
\item QSUP-MIX $\Rightarrow$ QSUP-PUR uses neither that $\rho$ is mixed nor that
  $p_{\hat\phi} > 0$; only $p_\phi > 0$, $p_{\hat\phi} \ge 0$ and $\rho_\phi \notin
  \{\rho_{\psi_1}, \rho_{\psi_2}\}$.
\end{itemize}

\paragraph{Problems found in the brief.} None that affect the theorem; the demonstration leaves
out some steps, which the formalization supplies.
\begin{itemize}
\item QSUP-PUR $\Rightarrow$ QSUP-MIX: the brief does not check that $\hat\phi \ne 0$ (needed to
  form $\rho_{\hat\phi}$), that $\rho_\phi \notin \{\rho_{\psi_1}, \rho_{\psi_2}\}$, or that $\rho$
  is mixed. All three follow from the linear independence of $\psi_1, \psi_2$ and
  $c_1, c_2 \ne 0$.
\item QSUP-MIX $\Rightarrow$ QSUP-PUR: ``Since $\rho$ is the mixture of two pure states, it is
  supported by a two dimensional subspace'' presupposes that $\psi_1$ and $\psi_2$ are linearly
  independent, which QSUP-MIX does not state, and the brief does not show that $c_1, c_2 \ne 0$.
  Both follow from $\rho_\phi \notin \{\rho_{\psi_1}, \rho_{\psi_2}\}$: if $\psi_2$ were a
  multiple of $\psi_1$, or $c_1$ or $c_2$ were $0$, then $\phi$ would be a multiple of $\psi_1$ or
  of $\psi_2$. The support claim itself is proved from the quadratic form: a vector orthogonal to
  $\psi_1$ and $\psi_2$ gives $\braket{x|\rho|x} = 0$, hence $\braket{\phi|x} = 0$.
\end{itemize}

\end{document}
